\noindent{\large\bfseries ملخص}\par\medskip

يقدّم هذا التقرير دراسة مقارنة لتنبّؤ أعطال المحامل الدوّارة اعتمادًا على البيانات الحقيقية لمنصّة PRONOSTIA (معهد FEMTO-ST، تحدّي IEEE PHM 2012). تتبع موثوقية مجموعة المحامل، المقدَّرة بطريقة الإمكان الأعظم مع مراعاة البيانات المبتورة، قانون فايبل ($\beta = 1{,}85$): B10 = 125 دقيقة وB50 = 346 دقيقة، مع مجالات ثقة واسعة. يُكتشف أوّل عيب بقاعدة الثلاثة انحرافات معيارية قبل العطل بـ 7,8 دقائق في المتوسّط الوسيط، ويُبنى مؤشّر للحالة الصحّية من طاقات حزم التحويل المويجي. تُقارَن سبع طرق لتقدير العمر المتبقّي: عملية فينر البايزية هي الأدقّ (جذر متوسّط مربّع الخطأ 16,4 دقيقة) وشبكة LSTM هي الأكثر حذرًا. وتؤكّد الاختبارات على بيانات محاكاة والاختبارات الإحصائية وتحليل قابلية التفسير صحّة المنهجية. وأخيرًا، تخفّض الصيانة التنبّؤية المعتمدة على نموذج فينر كلفة الساعة بنسبة 47 ٪ مقارنةً بالصيانة التصحيحية.

\textbf{الكلمات المفاتيح:} الموثوقية، قانون فايبل، الإمكان الأعظم، العمر المتبقّي، عملية فينر، تعلّم الآلة، الصيانة التنبّؤية.
